\SourceLeaf{405}%
wo \(c\) den Unterschied der Halbmesser des ersten und zweiten, und \(c'\) den Unterschied der Halbmesser des ersten und dritten Kreises bedeutet.

Wir haben folgende Gleichungen:
\[
\begin{aligned}
xx+yy&=zz\\
(x-a)^2+(y-b)^2&=(z-c)^2\\
(x-a')^2+(y-b')^2&=(z-c')^2
\end{aligned}
\]
Setzt man
\[
\begin{aligned}x&=z\cos\varphi\\y&=z\sin\varphi\end{aligned}
\]
so erhält man aus den vorigen Gleichungen
\[
\begin{aligned}
aa+bb-cc&=2az\cos\varphi+2bz\sin\varphi-2cz\\
a'a'+b'b'-c'c'&=2a'z\cos\varphi+2b'z\sin\varphi-2c'z
\end{aligned}
\]
Dividirt man die erste Gleichung mit \(a\cos\varphi+b\sin\varphi-c\), die zweite mit \(a'\cos\varphi+b'\sin\varphi-c'\), und zieht sie dann von einander ab, so erhält man
\[
\frac{a'a'+b'b'-c'c'}{a'\cos\varphi+b'\sin\varphi-c'}-\frac{aa+bb-cc}{a\cos\varphi+b\sin\varphi-c}=0
\]
oder wenn wir bezeichnen
\[
\begin{aligned}a'a'+b'b'-c'c'&=A'\\aa+bb-cc&=A\end{aligned}
\]
\[
(A'a-Aa')\cos\varphi+(A'b-Ab')\sin\varphi=A'c-Ac'
\]
Setzen wir nun
\[
\begin{aligned}A'a-Aa'&=R\cos M\\A'b-Ab'&=R\sin M\\A'c-Ac'&=N\end{aligned}
\]
so verwandelt sich unsere Gleichung in
\[
R\cos(\varphi-M)=N
\]
worin ausser \(\varphi\) alles gegeben ist. Man erhält daraus zwei Werthe für \(\varphi\), unter denen man nach der Art, wie der gesuchte Kreis berühren soll, zu wählen hat.
